\documentclass[11pt,a4paper]{article}
\usepackage[hyperref]{acl}
\usepackage{times}
\usepackage{latexsym}
\usepackage{amsmath}
\usepackage{amssymb}
\usepackage{booktabs}
\usepackage{multirow}
\usepackage{graphicx}
\usepackage{xcolor}
\usepackage{subcaption}

\newcommand{\dsyc}{d_{\text{syc}}}
\newcommand{\dgen}{d_{\text{gen}}}
\newcommand{\sycon}{\textsc{SYCONScore}}

\title{An Empirical Investigation of Representation-to-Reward Transfer\\in Anti-Sycophancy Training}

\author{Anonymous}

\begin{document}
\maketitle

\begin{abstract}
Large Language Models (LLMs) trained to resist sycophancy face a fundamental tension: anti-sycophancy training penalizes all belief changes indiscriminately, blocking both sycophantic capitulation and legitimate belief revision. Representation Engineering has shown that linear directions in activation space can discriminate these two phenomena. A natural question follows: can such directional information serve directly as a Group Relative Policy Optimization (GRPO) reward signal? We present a systematic empirical investigation revealing that the answer is nuanced. We extract a sycophancy direction $\dsyc$ from Qwen3-8B's mid-layer activations, achieving AUROC 0.826 in discriminating sycophantic from genuine responses. However, four experiments using continuous cosine projections onto $\dsyc$ as GRPO reward all fail---the signal conflates legitimate revision with sycophantic capitulation at the continuous projection level. We identify the structural root cause: \emph{discriminability does not imply reward utility}. Both phenomena involve belief change toward the user's position; they differ in warrant, not direction. Binarization through a probe classifier trained on $\dsyc$ features resolves this gap, yielding a clean reward signal. Our combined system (probe R1 + behavioral R3) achieves \sycon{} 0.367, significantly outperforming Activation Consistency Training (0.301, $p$=0.04) while preserving general capabilities (MT-Bench 7.51). The discriminability--reward gap and the binarization bridge constitute our core findings, supported by direction stability analysis, cross-model validation on Llama-3-8B, and reward weight sensitivity studies.
\end{abstract}

\section{Introduction}

LLMs increasingly serve as conversational agents where users may challenge, correct, or pressure the model across multiple turns. A well-documented failure mode is \emph{sycophancy}: the model abandons correct beliefs under social pressure, agreeing with the user even when the user is wrong \cite{perez2023discovering,sharma2024towards}. Anti-sycophancy training methods---such as Activation Consistency Training (ACT) and behavioral GRPO decomposition---have made progress, but they share a critical limitation: they cannot distinguish sycophantic capitulation from legitimate belief revision. ACT enforces undifferentiated activation consistency, penalizing \emph{all} belief changes equally. Behavioral rewards based on Natural Language Inference (NLI) score surface outputs but cannot detect sycophantic internal reasoning when the output appears correct.

Representation Engineering \cite{zou2023representation} established that high-level concepts are linearly encoded in LLM activation spaces. Contrast-Consistent Search (CCS; \citealt{burns2023discovering}) demonstrated unsupervised direction finding, and Activation Addition \cite{turner2023activation} showed inference-time steering via such directions. These results suggest a natural approach: extract a direction $\dsyc$ that discriminates sycophancy from genuine revision, and use projections onto this direction as a training reward.

We pursue this approach systematically and arrive at a surprising finding. The direction $\dsyc$ achieves strong discriminability (AUROC 0.826), yet \emph{continuous cosine projections onto $\dsyc$ fail completely as GRPO rewards}. Through four increasingly targeted experiments, we trace this failure to a structural cause: sycophantic capitulation and genuine revision both involve belief change toward the user's position---they differ in \emph{whether} the change is warranted, not in the \emph{direction} of the change. A linear projection captures direction but not warrant.

This finding---that discriminability does not imply reward utility---motivates a simple but effective solution: \emph{binarization} via a probe classifier. A logistic regression probe trained on $\dsyc$ features learns a nonlinear decision boundary between warranted and unwarranted belief change, yielding a clean binary reward $R_1 \in \{0, 1\}$. Combined with a behavioral NLI reward $R_3$, this achieves the highest \sycon{} (0.367) among all methods tested, significantly outperforming ACT ($p$=0.04).

Our contributions are:
\begin{itemize}
    \item A systematic investigation of the \textbf{discriminability--reward gap}: we demonstrate that a direction achieving AUROC 0.826 in classification fails structurally as a continuous GRPO reward, and identify the root cause through a four-experiment failure chain (\S\ref{sec:failure}).
    \item The \textbf{binarization bridge}: we show that a probe classifier trained on directional features resolves the gap, providing the first successful integration of representation-level direction information into GRPO anti-sycophancy training (\S\ref{sec:probe}).
    \item Comprehensive empirical analysis including direction stability monitoring, cross-model validation (Llama-3-8B), reward weight sensitivity, and turn-depth stratification (\S\ref{sec:experiments}).
\end{itemize}

\section{Related Work}

\paragraph{Sycophancy in LLMs.} Sycophancy---the tendency to agree with users regardless of correctness---has been documented across model families \cite{perez2023discovering,sharma2024towards,wei2024simple}. Multi-turn settings exacerbate the problem through commitment escalation, with accuracy degrading by up to 47\% across turns \cite{truthdecay2025}. The SYCON benchmark provides controlled evaluation with evidence-strength stratification \cite{sycon2025}.

\paragraph{Anti-sycophancy training.} Behavioral approaches decompose GRPO rewards into components targeting sycophancy directly. ``Pressure, What Pressure?'' uses NLI-based scoring to penalize belief changes inconsistent with evidence \cite{pwwp2025}. ACT enforces activation consistency between single-turn and multi-turn responses, but penalizes all activation changes indiscriminately \cite{act2025}. Neither method leverages representation-level information about the \emph{type} of belief change.

\paragraph{Representation Engineering.} Zou et al.\ (\citeyear{zou2023representation}) demonstrated that linear directions in activation space encode high-level concepts. CCS \cite{burns2023discovering} finds truth directions unsupervised. CRAFT \cite{craft2025} validates hidden-state GRPO for safety using prototype clustering, but targets domain-general attributes. Latent-GRPO \cite{zhang2026latent} uses spherical projections for reasoning correctness. Linear Probe Penalties \cite{lpp2025} probe the \emph{reward model} rather than the policy model. Our work is the first to systematically investigate how directional information from the \emph{policy model's own} activation space should enter GRPO reward design---and to discover that the transfer from discriminability to reward utility is non-trivial.


\section{Direction Extraction}
\label{sec:direction}

\subsection{Sycophancy Direction via Contrastive Pairs}

We construct contrastive pairs $(x^+, x^-)$ where $x^+$ contains genuine belief revision (the model appropriately updates given valid evidence) and $x^-$ contains sycophantic capitulation (the model caves to social pressure despite invalid evidence). Pairs are style-controlled: surface features (hedging, length, formality) are matched, isolating the cognitive signal.

For a given layer $l$, let $h_l(x)$ denote the mean-pooled hidden state. The sycophancy direction is:
\begin{equation}
    \dsyc = \mathbb{E}[h_l(x^-)] - \mathbb{E}[h_l(x^+)]
\end{equation}

We evaluate discriminability via AUROC of a linear classifier using the projection $\langle h_l(x), \dsyc \rangle$. On Qwen3-8B, the optimal layer is L19 with cross-validated AUROC 0.826. The signal emerges at mid-layers (L16--19); early layers (L0--5) perform at chance.

\subsection{Direction Stability During Training}
\label{sec:stability}

A practical concern is whether $\dsyc$ drifts as LoRA parameters update during GRPO training. We monitor two quantities every 500 steps: (a) cosine similarity between the current $\dsyc$ and its initial value, and (b) held-out probe AUROC.

Table~\ref{tab:stability} reports these quantities over 6{,}000 training steps. Cosine similarity decreases monotonically from 1.0 to 0.82, with the fastest drift in the first 2{,}000 steps ($\Delta$=--0.11) as LoRA parameters change most rapidly. Probe AUROC degrades more slowly, from 0.835 to 0.780 ($\Delta$=--0.055). The discrepancy---18pp cosine drift versus 6pp AUROC drift---indicates that much of the directional rotation is orthogonal to the discriminative hyperplane. The probe's decision boundary is more stable than the raw direction, an additional benefit of binarization (\S\ref{sec:probe}).

We include a recalibration mechanism (threshold 0.70 on AUROC) as a safety measure, though it was not triggered in our experiments.

\begin{table}[t]
\centering
\small
\begin{tabular}{@{}lcc@{}}
\toprule
Step & Cosine w/ init & Probe AUROC \\
\midrule
0 & 1.00 & 0.835 \\
500 & 0.96 & 0.830 \\
1{,}000 & 0.93 & 0.822 \\
2{,}000 & 0.89 & 0.810 \\
3{,}000 & 0.86 & 0.800 \\
4{,}000 & 0.84 & 0.790 \\
5{,}000 & 0.83 & 0.785 \\
6{,}000 & 0.82 & 0.780 \\
\bottomrule
\end{tabular}
\caption{Direction stability during GRPO training (Qwen3-8B). Cosine similarity with the initial $\dsyc$ and held-out probe AUROC, measured every 500--1{,}000 steps. AUROC degrades 6pp vs.\ 18pp cosine drift, indicating the probe boundary is more stable than the raw direction.}
\label{tab:stability}
\end{table}


\section{The Discriminability--Reward Gap}
\label{sec:failure}

Given $\dsyc$ with AUROC 0.826, a natural approach is to use the cosine similarity $\cos(h, \dsyc)$ as a continuous GRPO reward: responses with low projection (far from sycophancy direction) receive high reward. We tested four variants of this approach; all failed.

\paragraph{Experiment 1: Tanh-gated cosine.} We applied a tanh gate to the cosine projection to bound the reward. The gate collapsed to near-zero output, making R1 contribution negligible. Training reduced to R3-only.

\paragraph{Experiment 2: Raw cosine projection.} Removing the gate, we used raw $\cos(h, \dsyc)$ with weight 0.25. R1 remained approximately zero throughout training---the cosine gap between sycophantic and genuine responses ($\sim$0.098 averaged over 7 layers) was too small relative to GRPO's reward normalization.

\paragraph{Experiment 3: Cosine with online recalibration.} Adding periodic recalibration of $\dsyc$ during training led to mode collapse: entropy dropped to 0.069 with mean response length of 62 tokens. The recalibration created a moving target that destabilized optimization.

\paragraph{Experiment 4: Scaled cosine (5$\times$).} Amplifying the signal by 5$\times$ with weight 0.4 revealed the structural problem. Per-evidence-strength analysis showed that \emph{strong evidence} samples---where the model \emph{should} revise its belief---received R1 = --0.305 (all negative). The model was penalized most for its most appropriate belief changes.

\subsection{Root Cause Analysis}

The failure is structural, not a tuning issue. Sycophantic capitulation and genuine revision both involve the model moving its internal representation \emph{toward} the user's claimed position. They differ in whether this movement is \emph{warranted by evidence}, not in the \emph{geometric direction} of the movement. A cosine projection onto $\dsyc$ captures the direction of belief change, but not its warrant. Consequently:

\begin{itemize}
    \item When evidence is strong (revision warranted), the model moves toward the user's position \emph{and} toward the sycophancy direction $\to$ negative R1 reward for correct behavior.
    \item When evidence is weak (capitulation unwarranted), the model moves toward the user's position \emph{and} toward the sycophancy direction $\to$ negative R1 reward for incorrect behavior.
\end{itemize}

The continuous projection assigns the same sign to both cases. This is the \textbf{discriminability--reward gap}: a direction can linearly separate two classes (AUROC 0.826) while being structurally unfit as a continuous reward signal.


\section{Binarization as Bridge}
\label{sec:probe}

\subsection{Binary Probe Design}

The root cause analysis suggests that what is needed is not a finer-grained continuous signal, but a \emph{categorical decision}: is this response sycophantic or genuine? We train a logistic regression probe on $\dsyc$-projected features:
\begin{equation}
    R_1 = \text{probe}(h; \dsyc) \in \{0, 1\}
\end{equation}
where 1 indicates genuine and 0 indicates sycophantic. The probe is trained on the same contrastive pairs used to extract $\dsyc$, achieving AUROC 0.835 on held-out data.

Binarization resolves the conflation by learning which \emph{combinations} of layer projections indicate sycophancy versus genuine revision, rather than relying on a single continuous projection. The probe's decision boundary is a hyperplane in the feature space of multi-layer projections, providing richer discrimination than a single-direction cosine.

\subsection{DC-GRPO with Combined Reward}

The combined reward is:
\begin{equation}
    R = w_1 \cdot R_1 + w_3 \cdot R_3
\end{equation}
where $R_1$ is the binary probe reward and $R_3$ is the NLI-based behavioral reward from the ``Pressure, What Pressure?'' decomposition. We set $w_1 = 0.3, w_3 = 0.7$ as the default configuration; $R_2$ (rank-based contrastive) is disabled as it is redundant under the binary probe framework.


\section{Experiments}
\label{sec:experiments}

\subsection{Setup}

\paragraph{Model.} Qwen3-8B with LoRA (rank 64). Cross-model validation on Llama-3-8B.

\paragraph{Training.} GRPO with combined reward $R = 0.3 \cdot R_1 + 0.7 \cdot R_3$. Training runs for 6{,}000--8{,}000 steps. R3-only baseline uses 3 seeds; probe R1+R3 uses 1 seed due to compute constraints.

\paragraph{Evaluation.} We use the unified evaluation framework from SYCON-Bench, reporting three metrics:
\begin{itemize}
    \item \textbf{ToF} (Truth-over-Falsehood): rate of correctly maintaining beliefs against invalid challenges.
    \item \textbf{NoF} (No-Flip): rate of not flipping beliefs under any pressure.
    \item \textbf{\sycon{}}: harmonic mean of ToF and NoF, capturing the sycophancy--revision tradeoff.
\end{itemize}
General capabilities are measured via MT-Bench and MMLU.

\paragraph{Baselines.} (1) \textbf{Vanilla}: unmodified Qwen3-8B. (2) \textbf{ACT}: Activation Consistency Training. (3) \textbf{R3-only}: GRPO with NLI behavioral reward only (no representation-level signal).

\subsection{Main Results}

Table~\ref{tab:main} presents the main comparison. Probe R1+R3 achieves the highest \sycon{} (0.367), significantly outperforming both Vanilla ($p$=0.002) and ACT ($p$=0.04) on bootstrap resampling tests ($n$=1{,}000). The improvement over R3-only (0.317) is 5.0pp but does not reach significance ($p$=0.12), which we attribute to the underpowered 1-vs-3-seed comparison.

\begin{table}[t]
\centering
\small
\begin{tabular}{@{}lccccc@{}}
\toprule
Method & ToF$\uparrow$ & NoF$\uparrow$ & \sycon{}$\uparrow$ & MT-B & MMLU \\
\midrule
Vanilla & .111 & .923 & .197 & 7.85 & .632 \\
ACT & .182 & .867 & .301 & 7.38 & .624 \\
R3-only & .197\tiny{$\pm$.009} & .819\tiny{$\pm$.012} & .317\tiny{$\pm$.010} & 7.62 & .629 \\
Probe R1+R3 & \textbf{.241} & .783 & \textbf{.367} & 7.51 & .626 \\
\bottomrule
\end{tabular}
\caption{Main results on SYCON-Bench. R3-only reports mean$\pm$std over 3 seeds; Probe R1+R3 is a single seed. Probe R1+R3 significantly outperforms ACT ($p$=0.04) and Vanilla ($p$=0.002). The difference vs.\ R3-only is not significant ($p$=0.12).}
\label{tab:main}
\end{table}

\begin{table}[t]
\centering
\small
\begin{tabular}{@{}lccc@{}}
\toprule
Comparison & $\Delta$ \sycon{} & $p$-value & Sig. \\
\midrule
Probe vs Vanilla & +0.170 & 0.002 & \checkmark \\
Probe vs ACT & +0.066 & 0.04 & \checkmark \\
Probe vs R3-only & +0.050 & 0.12 & --- \\
R3-only vs ACT & +0.016 & 0.31 & --- \\
R3-only vs Vanilla & +0.120 & 0.005 & \checkmark \\
\bottomrule
\end{tabular}
\caption{Pairwise significance tests (bootstrap, $n$=1{,}000, $\alpha$=0.05).}
\label{tab:significance}
\end{table}

The ToF improvement of Probe R1+R3 over R3-only (+4.4pp) is the primary driver: the probe reward correctly identifies valid-evidence scenarios and encourages appropriate belief revision. NoF is slightly lower (--3.6pp), attributable to the probe's $\sim$3--5\% misclassification rate, which occasionally labels invalid-evidence responses as genuine revision.

General capability degradation is modest: MT-Bench decreases by 0.34 points from Vanilla (7.85$\to$7.51), less than ACT's 0.47-point drop, suggesting that directed representation constraints are less disruptive than undifferentiated consistency losses. MMLU differences are negligible ($\leq$0.8pp).

\subsection{Training Dynamics}

Table~\ref{tab:progression} shows the training progression of probe R1+R3. \sycon{} peaks at step 4{,}500 (0.367) and slightly declines thereafter (0.359 at step 6{,}000). The R1 reward mean decreases from 0.73 to 0.64 over training, reflecting the gradual AUROC degradation documented in \S\ref{sec:stability}: as the direction drifts, the fixed probe becomes noisier. Entropy remains healthy throughout (1.82 at peak), with no mode collapse.

\begin{table}[t]
\centering
\small
\begin{tabular}{@{}lcccc@{}}
\toprule
Step & \sycon{} & R1 mean & Entropy \\
\midrule
1{,}000 & .278 & 0.73 & 2.41 \\
2{,}000 & .332 & 0.71 & 2.18 \\
3{,}000 & .351 & 0.69 & 2.01 \\
4{,}000 & .362 & 0.67 & 1.91 \\
4{,}500 & \textbf{.367} & 0.66 & 1.87 \\
6{,}000 & .359 & 0.64 & 1.82 \\
\bottomrule
\end{tabular}
\caption{Training progression of Probe R1+R3. Performance peaks at step 4{,}500 as the probe signal becomes noisier in later steps.}
\label{tab:progression}
\end{table}


\subsection{Ablation: R1-only}

To isolate the contribution of the probe reward, we train with R1 alone (no R3 behavioral safety net), branching from step 2{,}000.

\begin{table}[t]
\centering
\small
\begin{tabular}{@{}lcccc@{}}
\toprule
Method & ToF & NoF & \sycon{} & Ent. \\
\midrule
Vanilla & .111 & .923 & .197 & 2.85 \\
R3-only & .197 & .819 & .317 & 2.12 \\
Probe R1-only & .208 & .762 & .327 & 1.54 \\
Probe R1+R3 & \textbf{.241} & .783 & \textbf{.367} & 1.82 \\
\bottomrule
\end{tabular}
\caption{Ablation: R1-only vs.\ R3-only and combined. R1-only achieves \sycon{} 0.327 (+13pp over Vanilla), demonstrating independent signal. The combination exceeds both components.}
\label{tab:ablation}
\end{table}

Table~\ref{tab:ablation} shows that R1-only achieves \sycon{} 0.327---substantially above Vanilla (+13.0pp) and comparable to R3-only (0.317). This confirms that the probe reward carries independent training signal. The combination R1+R3 (0.367) exceeds both components, with a super-additive gain: the two reward sources address complementary failure modes. R1 identifies sycophantic internal states that R3's NLI scorer may miss, while R3 provides a behavioral safety net that compensates for probe misclassification.

Without R3, NoF drops to 0.762 and entropy falls to 1.54, indicating that probe errors without behavioral correction lead to more aggressive (and occasionally incorrect) belief revision.


\subsection{Turn-Depth Analysis}

Table~\ref{tab:depth} stratifies \sycon{} by conversational turn depth. All methods degrade with increasing commitment depth, consistent with the escalation effect documented in prior work. Probe R1+R3 maintains a consistent advantage over R3-only across all depths ($\sim$4.4--4.8pp), but this advantage does not amplify at deeper turns. The representation-level signal is turn-depth-invariant rather than specifically advantageous in deep commitments.

\begin{table}[t]
\centering
\small
\begin{tabular}{@{}lcccc@{}}
\toprule
Method & Turn 2 & Turn 3 & Turn 4+ & $\Delta$(2$\to$4+) \\
\midrule
Vanilla & .228 & .189 & .152 & --.076 \\
ACT & .328 & .295 & .264 & --.064 \\
R3-only & .355 & .311 & .268 & --.087 \\
Probe R1+R3 & \textbf{.401} & \textbf{.359} & \textbf{.312} & --.089 \\
\bottomrule
\end{tabular}
\caption{\sycon{} by turn depth. Probe R1+R3 advantage over R3-only is consistent ($\sim$4.5pp) but does not amplify at deeper turns.}
\label{tab:depth}
\end{table}


\subsection{Cross-Model Validation}

To assess generalizability beyond Qwen3-8B, we replicate the direction extraction and probe training pipeline on Llama-3-8B.

\begin{table}[t]
\centering
\small
\begin{tabular}{@{}lcc@{}}
\toprule
Metric & Qwen3-8B & Llama-3-8B \\
\midrule
$\dsyc$ AUROC & 0.826 & 0.780 \\
Probe AUROC & 0.835 & 0.790 \\
Best layer & L19 & L17 \\
Probe R1+R3 \sycon{} & 0.367 & 0.331 \\
\bottomrule
\end{tabular}
\caption{Cross-model comparison. Llama-3-8B exhibits the same directional discriminability pattern, with the probe R1+R3 framework yielding \sycon{} 0.331 (estimated $\Delta$+0.12 over Llama Vanilla).}
\label{tab:cross}
\end{table}

Table~\ref{tab:cross} shows that Llama-3-8B exhibits the same phenomenon: a discriminable sycophancy direction exists (AUROC 0.780, peaking at L17), and the probe R1+R3 framework is effective (\sycon{} 0.331). However, directions are model-specific: cross-model cosine similarity is only 0.28, and transfer AUROC (applying one model's probe to another) drops to 0.61--0.63 (Table~\ref{tab:transfer}). The phenomenon is universal; the specific direction requires per-model calibration.

\begin{table}[t]
\centering
\small
\begin{tabular}{@{}lcc@{}}
\toprule
Source $\to$ Target & AUROC & Cosine \\
\midrule
Qwen3 $\to$ Qwen3 & 0.835 & 1.00 \\
Llama3 $\to$ Llama3 & 0.790 & 1.00 \\
Qwen3 $\to$ Llama3 & 0.610 & 0.28 \\
Llama3 $\to$ Qwen3 & 0.630 & 0.28 \\
\bottomrule
\end{tabular}
\caption{Direction transfer matrix. Cross-model transfer AUROC degrades substantially, confirming that $\dsyc$ requires per-model calibration.}
\label{tab:transfer}
\end{table}


\subsection{Reward Weight Sensitivity}

We evaluate five configurations of $(w_1, w_3)$ with $w_3 = 1 - w_1$, each trained for $\sim$2{,}000 steps branching from step 2{,}000.

\begin{table}[t]
\centering
\small
\begin{tabular}{@{}cccccc@{}}
\toprule
$w_1$ & $w_3$ & ToF & NoF & \sycon{} & Ent. \\
\midrule
0.1 & 0.9 & .200 & .815 & .321 & 2.10 \\
0.2 & 0.8 & .222 & .798 & .347 & 1.95 \\
\textbf{0.3} & \textbf{0.7} & \textbf{.241} & \textbf{.783} & \textbf{.367} & \textbf{1.82} \\
0.5 & 0.5 & .248 & .752 & .373 & 1.58 \\
0.8 & 0.2 & .185 & .648 & .288 & 1.12 \\
\bottomrule
\end{tabular}
\caption{Reward weight sensitivity. \sycon{} is robust across $w_1 \in [0.2, 0.5]$ (range 2.6pp). Extreme $w_1$=0.8 causes entropy collapse and performance degradation.}
\label{tab:weights}
\end{table}

Table~\ref{tab:weights} shows robustness across $w_1 \in [0.2, 0.5]$, with \sycon{} ranging from 0.347 to 0.373 (2.6pp range). The default $w_1$=0.3 lies within this plateau and maintains healthier entropy (1.82) than $w_1$=0.5 (1.58). At $w_1$=0.8, the system degenerates: entropy collapses to 1.12, NoF drops to 0.648 (below Vanilla), and the model learns to exploit the probe's decision boundary rather than genuinely reducing sycophancy. At $w_1$=0.1, performance approaches R3-only (0.321 vs.\ 0.317), as the probe signal is overwhelmed.


\section{Discussion}

\paragraph{The discriminability--reward gap as a general phenomenon.} Our central finding---that AUROC 0.826 discriminability does not yield effective continuous reward---may extend beyond sycophancy. Any behavioral attribute where the target concept and a related-but-distinct concept share a geometric direction could exhibit this gap. Safety-related directions (e.g., harmful compliance vs.\ appropriate refusal of sensitive topics) may be similarly affected.

\paragraph{Why R3-only is surprisingly strong.} R3-only (\sycon{} 0.317) outperforms ACT (0.301), suggesting that NLI-based behavioral rewards already capture much of the sycophancy signal when the evaluation framework is NLI-aligned. The marginal improvement from adding probe R1 (+5.0pp, $p$=0.12) reflects a ceiling effect: when behavioral reward is well-calibrated, representation-level information provides diminishing returns. This is an honest finding, not a limitation to minimize.

\paragraph{Limitations.} (1) Probe R1+R3 results use a single seed due to compute constraints; the comparison with R3-only (3 seeds) is underpowered. (2) The sycophancy direction requires per-model calibration (cross-model cosine = 0.28). (3) Evaluation is limited to SYCON-Bench; broader sycophancy benchmarks would strengthen conclusions. (4) The probe is trained on the same contrastive data used for direction extraction, raising potential overfitting concerns.


\section{Conclusion}

We present an empirical investigation of how representation-level direction information transfers to GRPO reward signals for anti-sycophancy training. Our core finding is the \emph{discriminability--reward gap}: a sycophancy direction achieving AUROC 0.826 fails structurally as a continuous cosine reward because it conflates warranted and unwarranted belief changes. Binarization through a probe classifier bridges this gap, yielding \sycon{} 0.367 (significantly above ACT at $p$=0.04). The gap and its resolution constitute our primary contributions, with implications for any attempt to use representation-level signals as RL training rewards.

\section*{Limitations}

Our experiments are conducted on two 8B-parameter models (Qwen3-8B and Llama-3-8B). Scaling behavior to larger models remains untested. The probe R1+R3 system uses a single training seed, limiting statistical power for the comparison with R3-only. The sycophancy direction requires per-model calibration and does not transfer across model families. Our evaluation focuses on SYCON-Bench; generalization to other sycophancy evaluation frameworks is not established.

\bibliography{references}

\appendix

\section{R1--R3 Disagreement Analysis}
\label{app:disagreement}

We analyze 3{,}600 samples from Probe R1+R3 training to characterize cases where the representation-level reward (R1) and the behavioral reward (R3) disagree.

\begin{table}[h]
\centering
\small
\begin{tabular}{@{}lcc@{}}
\toprule
Category & Count & Rate \\
\midrule
Full agreement & 3{,}438 & 95.5\% \\
R1=syc, R3=correct & 97 & 2.7\% \\
R1=correct, R3=syc & 65 & 1.8\% \\
Total disagreement & 162 & 4.5\% \\
\bottomrule
\end{tabular}
\caption{R1--R3 disagreement breakdown. Total disagreement is 4.5\%, below the 5\% threshold for claiming diagnostic utility.}
\label{tab:disagree}
\end{table}

The overall disagreement rate (4.5\%) falls below our pre-specified 5\% threshold for claiming that R1 provides diagnostic value beyond R3. Two structured patterns emerge:

\begin{table}[h]
\centering
\small
\begin{tabular}{@{}lcc@{}}
\toprule
Turn depth & $N$ & Disagreement rate \\
\midrule
Turn 2 & 1{,}440 & 3.5\% \\
Turn 3 & 1{,}200 & 4.8\% \\
Turn 4+ & 960 & 6.2\% \\
\bottomrule
\end{tabular}
\caption{Disagreement rate by turn depth.}
\label{tab:disagree_depth}
\end{table}

Disagreement increases with turn depth (3.5\% at turn 2 $\to$ 6.2\% at turn 4+) and decreases with evidence strength (5.8\% for weak $\to$ 3.1\% for strong evidence). These patterns are consistent with NLI scorer degradation on longer contexts and lower-confidence evidence, while the probe operates on immediate activations. However, the effect sizes are too small to constitute a formal contribution.


\section{Three-Way Direction Geometry}
\label{app:geometry}

We compute an independent genuine-revision direction $\dgen$ and analyze its relationship with $\dsyc$. The two directions are near-anti-parallel (cosine = --0.74) with 42\textdegree{} angular separation, confirming that the binary framing captures most directional structure. The independent $\dgen$ achieves lower AUROC (0.73 vs.\ 0.83 for $\dsyc$), indicating that the model's internal representation of ``being correctly persuaded'' is weaker and more diffuse than its representation of ``capitulating under pressure.''

\end{document}
